\newcommand{\algorithmOverview}{
\begin{algorithm}[H]
  \caption{The proposed ITR estimator accounting for competing risks.}
  \label{alg:Overview}
  \begin{algorithmic}
  \STATE \textbf{Input:} $Z_i, A_i, \tol$
  \STATE \textbf{Output:} $\dopt$
  \STATE %\ref{alg:RSF}
  \STATE For each individual $i$, given covariates $Z_i$ and small $0<\tol<1$, consider feasible treatments $a_i \in \psi(z_i) 
  \subset \mathcal{A}$ (we drop $i$ for simplicity):
  \begin{enumerate}
    \item Obtain $\survfunchat(\cdot|Z,A)$ via generalized random survival forest \citep{cho_multi-stage_2023}
    \item Obtain $\hat{F}_j(\cdot|Z,A)$ via random cumulative incidence forest (Algorithm \ref{alg:RCIF})
    \item Let optimal treatments for Phase 1 (overall survival) be $\valueonemaxTRUE = \argmax_{a \in \psi(z)}\valueone$
    \item Obtain overall survival treatment set $A^s = \{a \in \psi(Z): \valueone \ge (1- \tol)\valueonemaxTRUE\}$
        \begin{itemize}
            \item If $A^s$ is singleton, $A^s$ is the optimal treatment: $\dopt(Z) = A^s$.
            \item If $\#A^s > 1$, optimal treatment $\dopt(Z) = \argmin_{a \in \textit{A}^s} \valuetwo$.
        \end{itemize} 
    \end{enumerate}
  \STATE In other words, if $a_i \neq A^s_i$ for individual $i$, then individual $i$ stopped at either 
  \begin{itemize*}%[label={Phase \arabic*:}]
      \item \hspace{1em} Phase 1: $\valueone < (1-\tol)\valueonemaxTRUE$ or
      \item \hspace{1em} Phase 2: $\valueone \ge (1-\tol) \valueonemaxTRUE$ but $V_2(a = \dopt, Z) < V_2(a = a_i, Z)$
  \end{itemize*}\\
  and $A_i^s$ may have $\#A_i^s>1$.
\end{algorithmic}
\end{algorithm}
}

\newcommand{\algorithmRSF}{
\begin{algorithm}[H]
  \caption{RSF}
  \label{alg:RSF}
  \begin{algorithmic}
  \STATE \textbf{Input:} $Z_i, A_i$
  \STATE \textbf{Output:} $A_i^*$
  \STATE %\ref{alg:RSF}
  \STATE For each individual $i$, given covariates $Z_i$ and small $0<\tol<1$, consider feasible treatments $a_i \in \psi(z_i) \subset \mathcal{A}$:
  \STATE
  \STATE hi
  \STATE
\end{algorithmic}
\end{algorithm}
}

\newcommand{\algorithmgRCIF}{%

\begin{algorithm}[H]
  \caption{the random cumulative incidence forest estimator}
  \label{alg:RCIF}
  \begin{algorithmic}
  \STATE Result: random cause-j cumulative incidence forest $\hat{F}_j(\cdot|Z,A)$
  \STATE Parameters: the number of trees $n_{\text{tree}}$, the number of minimum terminal node size $n_{\text{min}}$
  \STATE \textbf{Input:} $\{(X_i, \delta_{ji}, Z_i, A_i)\}_{i=1}^n$
  \STATE \textbf{Output:} $\hat{F}_j(\cdot|Z,A=a)$
  \STATE 
  \STATE for \underline{treatment arm} $a \in \mathcal{A}$ \textbf{do} \\
  \STATE \quad for \underline{tree $w = 1, 2, \ldots, n_{\text{tree}}$} \textbf{do} \\
  \STATE \quad \quad Partition the feature space with the input data $\{(I(t\leq X_i), \delta_{ji}, Z_i)\}_{i:a_i=a}$ via Gray's test \cref{grays_test} or MD \citep{cho_multi-stage_2023} until node size $< 2n_{\text{min}}$ \\
  \STATE \quad \quad Obtain the terminal node cumulative incidence probability $\hat{F}^a_{j,\{w\}}(t | Z)$ via Aalen-Johanson estimator \cref{AJ_estimator} \\
  \STATE \quad Obtain the random cumulative incidence forest estimator \\
  \STATE \quad $\hat{F}_j(\cdot | Z, A = a) = \frac{1}{n_{\text{tree}}} \sum_{w=1}^{n_\text{tree}} \hat{F}_{j,\{w\}}^{a}(t | Z)$ \\

\end{algorithmic}
\end{algorithm}
}